% !TEX TS-program = pdflatex
% نسخه فارسی گزارش پروژه دکترای RL
\documentclass[12pt]{article}
\usepackage[margin=1in]{geometry}
\usepackage{amsmath,amssymb}
\usepackage{graphicx}
\usepackage{booktabs}
\usepackage{hyperref}
\usepackage{enumitem}
\usepackage{xcolor}
\hypersetup{colorlinks=true,linkcolor=blue,citecolor=blue,urlcolor=blue}
\setlength{\parskip}{6pt}
\setlength{\parindent}{0pt}

\title{چارچوب‌های نظری و تجربی یادگیری تقویتی عمیق \\ نسخه فارسی گزارش جامع}
\author{مخزن درس یادگیری تقویتی عمیق}
\date{\today}

\begin{document}
\maketitle
\tableofcontents

\section{هدف گزارش}
این سند نسخه فارسی گزارش پروژه است و تمام زنجیره‌های پنج‌گانه‌ی مقاله‌ها (مبتنی بر ارزش، گرادیان سیاست، بازیگر--ناظر، مبتنی بر مدل، و چندعامله) را پوشش می‌دهد. هر بخش شامل توضیح نظری کوتاه، تنظیمات آزمایش، و جدول نتایج قابل پرشدن است.

\section{چگونه بازتولید کنیم}
\begin{itemize}[leftmargin=*]
    \item محیط مجازی: \verb|python3.10 -m venv .venv_runs| سپس \verb|source .venv_runs/bin/activate|.
    \item نصب وابستگی‌ها: \verb|pip install -r requirements.txt|
    \item اجرای یک آزمایش: \verb|python scripts/run_experiment.py --chain value --algo dqn --env CartPole-v1 --steps 20000 --seed 0|
    \item اجرای کل سوئیت: \verb|python scripts/run_suite.py --suite configs/suites/fast_5chain.yaml --mode full|
    \item ساخت نمودار و جدول: \verb|python scripts/generate_plots.py| سپس \verb|python scripts/build_report_tables.py|
    \item خروجی‌ها در \verb|outputs/runs| و \verb|outputs/suite_reports| ذخیره می‌شوند.
\end{itemize}

\section{زنجیره ۱: روش‌های مبتنی بر ارزش (DQN، Double، Dueling)}
\textbf{ایده اصلی:} بازپخش تجربه، شبکه هدف، کاهش سوگیری بیش‌برآورد با Double، و نمایش بهتر با Dueling.\\
\textbf{تنظیمات پیشنهادی برای CartPole:} نرخ یادگیری $10^{-3}$، \texttt{train\_freq}=1، اندازه بافر 50k، اپسیلون نهایی 0.02، \texttt{total\_steps}=100k.\\
\IfFileExists{outputs/suite_reports/figures/value_dqn_CartPole-v1.png}{
\begin{figure}[h]
    \centering
    \includegraphics[width=0.82\linewidth]{outputs/suite_reports/figures/value_dqn_CartPole-v1.png}
    \caption{نمودار یادگیری زنجیره ارزش (تولید خودکار).}
\end{figure}
}{\textit{نمودار خودکار این بخش هنوز تولید نشده است.}}
\textbf{جدول نتایج (پر شود):}
\begin{table}[h]
    \centering
    \begin{tabular}{lccc}
    \toprule
    الگوریتم & گام‌ها & بازده میانگین & انحراف معیار \\
    \midrule
    DQN & & & \\
    Double DQN & & & \\
    Dueling DQN & & & \\
    QR-DQN (اختیاری) & & & \\
    \bottomrule
    \end{tabular}
\end{table}

\section{زنجیره ۲: گرادیان سیاست و ناحیه اعتماد (REINFORCE، PPO، TRPO، CPO-lite)}
توضیح کوتاه: REINFORCE با پایه، TRPO با قید KL، PPO با برش نسبت احتمال، نسخه ساده CPO با جریمه لاگرانژی هزینه.\\
\IfFileExists{outputs/suite_reports/figures/policy_ppo_CartPole-v1.png}{
\begin{figure}[h]
    \centering
    \includegraphics[width=0.82\linewidth]{outputs/suite_reports/figures/policy_ppo_CartPole-v1.png}
    \caption{نمودار یادگیری روش‌های سیاست‌محور (تولید خودکار).}
\end{figure}
}{}
\textbf{جدول نتایج:}
\begin{table}[h]
    \centering
    \begin{tabular}{lcccc}
    \toprule
    الگوریتم & محیط & گام‌ها & بازده میانگین & هزینه میانگین \\
    \midrule
    REINFORCE & & & & \\
    PPO & & & & \\
    TRPO & & & & \\
    CPO-lite & & & & \\
    \bottomrule
    \end{tabular}
\end{table}

\section{زنجیره ۳: بازیگر--ناظر و آنتروپی بیشینه (A3C/A2C، SAC)}
توضیح کوتاه: A2C به‌عنوان جایگزین سبک A3C؛ SAC با آنتروپی خودکار برای کنترل پیوسته.\\
\IfFileExists{outputs/suite_reports/figures/actor_critic_sac_Pendulum-v1.png}{
\begin{figure}[h]
    \centering
    \includegraphics[width=0.82\linewidth]{outputs/suite_reports/figures/actor_critic_sac_Pendulum-v1.png}
    \caption{نمودار یادگیری بازیگر--ناظر (تولید خودکار).}
\end{figure}
}{}
\textbf{جدول نتایج:}
\begin{table}[h]
    \centering
    \begin{tabular}{lccc}
    \toprule
    الگوریتم & گام‌ها & بازده میانگین & انحراف معیار \\
    \midrule
    A2C (A3C proxy) & & & \\
    SAC & & & \\
    \bottomrule
    \end{tabular}
\end{table}

\section{زنجیره ۴: مبتنی بر مدل (Dyna-Q، ME-TRPO)}
توضیح کوتاه: Dyna-Q با به‌روزرسانی‌های برنامه‌ریزی؛ ME-TRPO با مدل تجمیع چندگانه برای کاهش سوگیری مدل.\\
\IfFileExists{outputs/suite_reports/figures/model_based_mbpo_lite_Pendulum-v1.png}{
\begin{figure}[h]
    \centering
    \includegraphics[width=0.82\linewidth]{outputs/suite_reports/figures/model_based_mbpo_lite_Pendulum-v1.png}
    \caption{نمودار یادگیری زنجیره مدل‌محور (تولید خودکار).}
\end{figure}
}{}
\textbf{جدول نتایج:}
\begin{table}[h]
    \centering
    \begin{tabular}{lccc}
    \toprule
    الگوریتم & محیط & معیار & مقدار \\
    \midrule
    Dyna-Q & CliffWalking-v0 & میانگین 100 اپیزود آخر & \\
    ME-TRPO & Pendulum-v1 & بازده ارزیابی & \\
    \bottomrule
    \end{tabular}
\end{table}

\section{زنجیره ۵: چندعامله (QMIX، MADDPG)}
توضیح کوتاه: QMIX با تابع اختلاط یکنوا، MADDPG با منتقد متمرکز و بازیگر غیرمتمرکز.\\
\IfFileExists{outputs/suite_reports/figures/marl_qmix_lite_simple_spread_v3.png}{
\begin{figure}[h]
    \centering
    \includegraphics[width=0.82\linewidth]{outputs/suite_reports/figures/marl_qmix_lite_simple_spread_v3.png}
    \caption{نمودار یادگیری چندعامله (تولید خودکار).}
\end{figure}
}{}
\textbf{جدول نتایج:}
\begin{table}[h]
    \centering
    \begin{tabular}{lcc}
    \toprule
    الگوریتم & تکرار & پاداش اپیزود میانگین \\
    \midrule
    QMIX & & \\
    MADDPG & & \\
    \bottomrule
    \end{tabular}
\end{table}

\section{جمع‌بندی}
پس از اجرای سوئیت، این فایل را با \texttt{latexmk -pdf report\_fa.tex} کامپایل کنید. برای جمع‌بندی، تفاوت پایداری، کارایی نمونه، ایمنی و مقیاس‌پذیری بین پنج زنجیره را بر اساس جدول‌ها و نمودارهای تولیدی تحلیل کنید.

\end{document}
